\documentclass[10pt]{article}
\usepackage{icmlko}

\icmlmeta{IP 파운데이션 월드모델}{257}{두 번째 문이 열렸다}{같은 절차가 애니에서도 통한다 --- 판 0.4623 $\to$ 0.4687, 애니 혼자 $+$0.042}

\begin{document}

\twocolumn[
\icmlheader

\begin{abstract}
노트 255가 웹툰 태그 \emph{내용}으로 판을 올렸고, 노트 256이 세계애니에서
같은 절차를 돌렸다가 떨어뜨리며 채택 검사에 \textbf{겹침}을 넣었다. 세
번째 시도가 애니다 --- 몫 23.2\% 로 웹툰 다음이고 \texttt{tags} 156종에
\texttt{genres} · \texttt{medium} 까지 있다. \textbf{한 번 미끄러졌다}:
웹툰 경계로 시간 조각을 자르니 다섯 중 둘이 60건 미만이라 검사가 통과할 수
없었다. 애니 학습은 2017--2025에 몰려 있어 경계가 다르다. \textbf{조각
경계는 도메인마다 제 연도 분위로} 잡아야 한다. 고쳐서 다시 재니 성분 둘이
통과한다 --- 연도 통제 $+$0.248(겹침 0.401) · $+$0.170(겹침
\textbf{0.184}), 조각 부호 둘 다 5/5. 애니의 \texttt{target\_breadth} 는
\textbf{연령 등급}이라 태그 내용과 다른 물건이고 \texttt{goods\_scale} 은
아예 비어 있어, 노트 256이 세계애니에서 만난 ``흐린 사본''이 아니다.
\textbf{판이 0.4623에서 0.4687로 올랐다}(F23 $+$0.0063 $t{=}1.84$, F21
$+$0.0125 $t{=}3.37$). 애니가 0.4868 $\to$ \textbf{0.5291}($+$0.0423,
$t{=}3.01$)이다. 값은 공짜가 아니다 --- \textbf{웹툰이 $-$0.0075}($t{=}-2.69$)
내고 나머지 일곱은 잡음이다. 웹툰만/애니만/둘 다를 따로 재면 $+$0.0054 ·
$+$0.0080 대 $+$0.0118 로 교차항이 $-$0.0016 이라, \textbf{노트 250의
``전용 이름은 가법''이 세 번째로 확인된다.} 두 노트 만에 판이 0.4569에서
0.4687로 왔다 --- 노트 252가 잰 유보 안 CV 천장(0.4697)에 \textbf{0.001}
남았다.
\end{abstract}
]

\section{세 번째 도메인}

절차는 노트 255가 세웠고 노트 256이 검사 하나를 더 붙였다. 남은 큰 도메인은
애니(몫 23.2\%, $\rho$ 0.487)다. 레코드를 열어 보니 \texttt{tags} 156종에
\texttt{genres} 18종 · \texttt{medium} 4종까지 있다 --- 웹툰보다 재료가 많다.

\section{조각 경계에서 미끄러졌다}

처음 돌렸을 때 통과한 성분이 \textbf{하나도 없었다}. 표를 보니 이상했다 ---
성분 2가 연도 통제 $+$0.248 인데 조각 부호가 ``3/3''이다. 다섯 조각을
쟀는데 셋만 세어졌다.

웹툰 경계(2016 · 2019 · 2021 · 2022 · 2023)를 그대로 썼기 때문이다. 애니
학습은 \textbf{2017--2025에 몰려 있어} 앞 두 조각이 60건 미만으로 비었다.
검사가 ``5/5''를 요구하는데 애초에 다섯이 안 만들어졌다.

\begin{center}\small
\begin{tabular}{lrrrrr}
\hline
애니 학습 분위 & 2017.3 & 2018.9 & 2020.3 & 2021.7 & 2023.0 \\
조각별 건수 & 293 & 292 & 295 & 292 & 295 \\
\hline
\end{tabular}
\end{center}

\textbf{조각 경계는 도메인마다 제 연도 분위로 잡는다.} 노트 187의 값싼
검사가 또 걸렸다 --- 결과가 ``전부 탈락''으로 깨끗하게 나오면 검사가 안
닿은 것을 의심할 것.

\section{둘이 통과한다}

\begin{center}\small
\begin{tabular}{lrrrl}
\hline
성분 & 연도 통제 & 조각 부호 & 겹침 & 상대 \\
\hline
0 & $+$0.069 & 3/5 & 0.106 & \\
1 & $+$0.009 & 3/5 & 0.142 & \\
\textbf{2} & $\mathbf{+0.248}$ & \textbf{5/5} & 0.401 &
  \texttt{target\_breadth} \\
\textbf{3} & $\mathbf{+0.170}$ & \textbf{5/5} & \textbf{0.184} &
  \texttt{venue\_prominence} \\
4 & $-$0.103 & 3/5 & 0.207 & \\
5 & $-$0.093 & 4/5 & 0.250 & \\
\hline
\end{tabular}
\end{center}

겹침이 0.401 · 0.184 로 웹툰(0.356 · 0.389)과 비슷하고 세계애니의 0.621
과는 멀다. 게다가 겹치는 상대가 다르다 --- 애니의
\texttt{target\_breadth} 는 \texttt{scale01(age, 0, 19)} 즉 \textbf{연령
등급}이고 \texttt{goods\_scale} 은 ``대응물 없음''으로 비어 있다. 노트 256이
만난 것(포맷 $=$ 총 분량)처럼 같은 물건을 두 번 말하는 구조가 아니다.

\section{판이 올랐다}

\begin{figure}[t]
\centering
\includegraphics[width=\columnwidth]{figs/board.pdf}
\caption{왼쪽 --- 판 사다리. 오른쪽 --- 애니 축을 더했을 때 도메인별 변화.
회색은 $|t| < 2$ 다.}
\end{figure}

\begin{center}\small
\begin{tabular}{lrr}
\hline
 & 판 & 짝 차 \\
\hline
노트 242 (17축) & 0.4569 & --- \\
$+$ 웹툰 태그 (19축) & 0.4623 & $+$0.0054 \\
$+$ 애니 태그만 (19축) & 0.4649 & $+$0.0080 \\
\textbf{둘 다 (21축)} & \textbf{0.4687} & $+$0.0118 \\
\hline
\end{tabular}
\end{center}

$+$0.0054 와 $+$0.0080 의 합이 $+$0.0134 인데 둘 다는 $+$0.0118 이다 ---
\textbf{교차항 $-$0.0016.} 노트 250이 ``전용 이름은 가법''(교차항/부분합
0.18)이라 한 것이 세 번째로 확인된다. 노트 249 · 251에서 공유 이름이
1.01 이었던 것과 대비된다.

애니 축만 놓고 보면 F21 $+$0.0125($t{=}3.37$) · F18 $+$0.0038($t{=}1.04$) ·
F23 $+$0.0063($t{=}1.84$)이다. \textbf{이번엔 능형이 얻는다} --- 노트 255의
웹툰 축은 나무가 얻었다(F18 $+$0.0053, F21 $-$0.0020). 같은 절차로 만든
축이 형식에 따라 다른 곳에서 값을 낸다.

\section{공짜가 아니다}

\begin{center}\small
\begin{tabular}{lrrr}
\hline
도메인 & 차 & 짝 SE & $t$ \\
\hline
\textbf{애니} & $\mathbf{+0.0423}$ & 0.0141 & \textbf{3.01} \\
아이돌 & $-$0.0377 & 0.0552 & $-$0.68 \\
\textbf{웹툰} & $\mathbf{-0.0075}$ & 0.0028 & $\mathbf{-2.69}$ \\
게임 & $-$0.0063 & 0.0046 & $-$1.36 \\
나머지 다섯 & $|{\cdot}| \le 0.005$ & & $|t| \le 1.3$ \\
\hline
\end{tabular}
\end{center}

애니가 $+$0.042 를 얻고 \textbf{웹툰이 $-$0.0075 를 낸다}($t{=}-2.69$).
아이돌의 $-$0.038 은 짝 SE 가 0.055 라 잡음이다(노트 248의 규약대로 읽었다).
\texttt{rank.spread} 로 보면 진짜 총량 0.0119 에 진짜 순 $+$0.0078 로
상쇄가 1.5배 --- 노트 247의 19배와 달리 대부분이 남는다.

웹툰이 왜 내는지는 아직 안 갈랐다. 전용 이름이라 계수는 안 나누지만
나무 문턱은 나눈다(노트 251). 다음 노트의 첫 줄이다.

\section{천장이 0.001 남았다}

\begin{center}\small
\begin{tabular}{lr}
\hline
 & 판 \\
\hline
노트 242 & 0.4569 \\
\textbf{지금} & \textbf{0.4687} \\
노트 252가 잰 유보 안 CV 천장 & 0.4697 \\
\hline
\end{tabular}
\end{center}

노트 252는 ``축 열일곱의 여지가 $+$0.013 뿐''이라고 적었고, 두 노트 만에
그 여지의 \textbf{90\%를 썼다.} 다만 그 천장은 \emph{열일곱 축}의 천장이고,
지금은 축이 스물하나다 --- 천장 자체를 다시 재야 한다. 그것이 이 노트가
남기는 제일 큰 숙제다.

\paragraph{남은 것}
\begin{center}\small
\begin{tabular}{ll}
\hline
갈래 & 왜 \\
\hline
천장 다시 재기 & 21축의 유보 안 CV 는 얼마인가 \\
웹툰 $-$0.0075 & 전용 이름인데 왜 내나. 나무 문턱? \\
모바일 · 게임 & \texttt{genres} 36 · 12종. 몫 25.5\% \\
펀딩 \texttt{category} & 20종, 몫 3.1\% \\
\hline
\end{tabular}
\end{center}

\paragraph{규약} 시간 조각 검사는 \textbf{도메인마다 제 연도 분위로} 자른다.
한 도메인의 경계를 다른 도메인에 그대로 쓰면 조각이 비고, 검사는 ``전부
탈락''이라는 깨끗한 거짓말을 한다.

\end{document}
